\subsection{nml\_chemistry}
\label{SEC:nml_chemistry}
\inp*[nmlChemistry]{}[nml\_chemistry][]

This \hyperref[SEC:User_Input]{namelist block}
is used to specify control of chemically reacting flows.

\begin{nmlblocklst}[
,label={LST:nml_chemistry}
,caption={Namelist nml\_chemistry.}
]
&nml_chemistry
  Disable_Chemical_Reactions       = F
  Semi_Implicit_Chemistry          = 0
  Ignition_Method                  = 0
  Ignition_Start_Timestep          = 0
  Ignition_Duration                = 1
  Ignition_Prod2Reac_L2_Ratio_Goal = -1.0
  Ignition_Max_Temperature         = 1600.0         ! [K]
  Ignition_Spark_Cap               = 1.0e+12        ! [N/m^2]
  Ignition_Source_Magnitude        = 20.0
  Ignition_Source_Radius           = 0.0            ! [m]
  Ignition_Source_Location         = [0.0,0.0,0.0]  ! [m]
/
\end{nmlblocklst}

%==========================================================================

\begin{nmldescitem}
Disable_Chemical_Reactions = F
\end{nmldescitem}
\inp*[nmlChemistry::DisableChemicalReactions]{}[Disable\_Chemical\_Reactions][]

\begin{nmldescbody}
If running a multi-component gas mixture, setting this to true will disable
all chemical reactions, \thatis all reaction source terms are set to zero,
resulting in a simulation of an inert gas mixture.
 
When reactions are disabled, all non-existent chemical species will be
dropped from the simulation to reduce computational cost. For example, if
running the 7-species ($\mathrm{H_2}$, $\mathrm{O_2}$, $\mathrm{H_2 O}$,
OH, H, O, $\mathrm{N_2}$)\slash 7-reaction mechanism for a hydrogen\slash
air mixture with reactions disabled and the flow initialized to a
mixture of $\mathrm{H_2}$, $\mathrm{O_2}$, and $\mathrm{N_2}$,
then the product species $\mathrm{H_2 O}$, OH, H, and O will be dropped
completely from the simulation.
%
If a simulation is run with reactions disabled, any dropped species
will not be written to the restart files. However, these non-reacting
restart files can still be used to restart a simulation with reactions
enabled and the product species will be added to the simulation with an
initial mass fraction of zero.
\end{nmldescbody}

%==========================================================================

\begin{nmldescitem}
Semi_Implicit_Chemistry = 0
\end{nmldescitem}
\inp*[nmlChemistry::SemiImplicitChemistry]{}[Semi\_Implicit\_Chemistry][]

\begin{nmldescbody}
If running a multi-component gas mixture with reactions, this input controls
the numerical treatment of the reaction source terms:

\begin{tabular}{ @{} r @{ = } p{0.85\linewidth} @{} }
   0 & Standard explicit treatment of the reaction source terms, which is
       likely to require the smallest time step.\\
   1 & Point implicit treatment of the reaction source terms for the species
       transport equations, using an analytical evaluation of the Jacobian
       matrix.\\
   2 & Point implicit treatment of the reaction source terms for the species
       transport equations, using a numerical approximation of the Jacobian
       matrix.\\
  -1 & Use special hard-coded source terms if using the 7-species \slash
       7-reaction hydrogen\slash air reaction mechanism.
       Otherwise, the same as option 0.
\end{tabular}

If the reaction terms require a smaller time step than what would be
required without reactions, then using a point implicit treatment of the
reaction source terms may (or may not) allow for using a time step closer
to the non-reacting case.
\vspace{3pt}
\end{nmldescbody}

%==========================================================================

\FINALFIX{\newpage}
\begin{nmldescitem}
Ignition_Method = 0
\end{nmldescitem}
\inp*[nmlChemistry::IgnitionMethod]{}[Ignition\_Method][]

\begin{nmldescbody}
Specifies the ignition method to use for speeding up the startup of chemical
reactions within the solution.

\begin{tabular}{ @{} r @{ = } p{0.85\linewidth} @{} }
   0 & Ignition process is disabled.\\
   1 & A source term of $\rho c_p \Delta T$ is added to the energy equation each
       Runge-Kutta stage, where $\Delta T = 
       \inp[nmlChemistry::IgnitionSourceMagnitude]{Ignition\_Source\_Magnitude}
       / [\#\ \textsf{RK Stages}]$.\\
   2 & A percentage of the reference total energy, $(\rho e_t)_{\infty}$, is
       added to the energy equation as a source term. The input
       \inptext{Ignition\_Source\_Magnitude} specifies the percentage to use, so
       the source term will be
       $\inptext{Ignition\_Source\_Magnitude}\times(\rho e_t)_{\infty}$.\\
   3 & The temperature value specified with
       \inptext{Ignition\_Source\_Magnitude} is used instead of the actual
       local temperature for the solution point when computing the species
       reaction source terms.
\end{tabular}
\end{nmldescbody}

%==========================================================================

\begin{nmldescitem}
Ignition_Start_Timestep = 0
\end{nmldescitem}
\inp*[nmlChemistry::IgnitionStartTimestep]{}[Ignition\_Start\_Timestep][]

\begin{nmldescbody}
Iteration number to begin using the ignition source terms for chemistry
reactions.
\end{nmldescbody}

%==========================================================================

\begin{nmldescitem}
Ignition_Duration = 1
\end{nmldescitem}
\inp*[nmlChemistry::IgnitionDuration]{}[Ignition\_Duration][]

\begin{nmldescbody}
Specifies the number of time steps for the ignition source term to remain
active once either the temperature at any solution point reaches
\inp[nmlChemistry::IgnitionMaxTemperature]{Ignition\_Max\_Temperature} or the
target goal for the overall ratio of L2(products)\slash L2(reactants),
\inp[nmlChemistry::IgnitionProd2ReacL2RatioGoal]{Ignition\_Prod2Reac\_L2\_Ratio\_Goal}, 
within the global solution is achieved. 
%
The ignition source terms turn off after this number of time steps are
completed.
\end{nmldescbody}

%==========================================================================

\begin{nmldescitem}
Ignition_Prod2Reac_L2_Ratio_Goal = -1.0
\end{nmldescitem}
\inp*[nmlChemistry::IgnitionProd2ReacL2RatioGoal]{}[Ignition\_Prod2Reac\_L2\_Ratio\_Goal][]

\begin{nmldescbody}
Target goal for the ratio of L2(products)\slash L2(reactants) of the simulation
to signify that reactions are occurring. Once this goal is met, the ignition
source terms will continue to be active for the number of time steps specified
by \inp[nmlChemistry::IgnitionDuration]{Ignition\_Duration}, after which, the
ignition source terms will be disabled.
\end{nmldescbody}

%==========================================================================

\begin{nmldescitem}
Ignition_Max_Temperature = 1600.0  ! [K]
\end{nmldescitem}
\inp*[nmlChemistry::IgnitionMaxTemperature]{}[Ignition\_Max\_Temperature][]

\begin{nmldescbody}
Specifies a value for temperature (in Kelvin) that marks when to begin turning
off the ignition source terms. Once this temperature has been reached, the
ignition source terms will continue to be active for the number of time steps
specified by \inp[nmlChemistry::IgnitionDuration]{Ignition\_Duration}, after
which, the ignition source terms will be disabled.
%
Any solution point that reaches this temperature while the ignition source terms
are still active will have its ignition source term set to zero.
%
When using $\inp[nmlChemistry::IgnitionMethod]{Ignition\_Method}=3$, this means
the actual local temperature at the solution point will be used to compute the
reaction source terms instead of the reaction temperature specified by
\inp[nmlChemistry::IgnitionSourceMagnitude]{Ignition\_Source\_Magnitude}.
\end{nmldescbody}

%==========================================================================

\begin{nmldescitem}
Ignition_Spark_Cap = 1.0e+12  ! [N/m^2]
\end{nmldescitem}
\inp*[nmlChemistry::IgnitionSparkCap]{}[Ignition\_Spark\_Cap][]

\begin{nmldescbody}
This specifies a maximum limit that can be used for the energy source term
(in $\mathrm{N}/\mathrm{m}^2$).
%
%\FIXME{
%Are these the correct units for an energy source term? YES.
%\inptext{Ignition\_Source\_Magnitude} uses Kelvin for ignition source term.
%When IgnSrcMag is Kelvin, it is multiplied by $\rho c_p$ to create the source
%term which is then compared to this value.
%}
\end{nmldescbody}

%==========================================================================

\begin{nmldescitem}
Ignition_Source_Magnitude = 20.0
\end{nmldescitem}
\inp*[nmlChemistry::IgnitionSourceMagnitude]{}[Ignition\_Source\_Magnitude][]

\begin{nmldescbody}
Value to use for the ignition source term. The value specified depends on the
ignition method specified by
\inp[nmlChemistry::IgnitionMethod]{Ignition\_Method}.

\begin{InputList}
  \item[$\inptext{Ignition\_Method}=1$]
    This input represents a temperature value (units of \unitTemp{}) that is
    used to compute a source term that is added
    to the energy equation each Runge-Kutta stage. The energy source term is
    evaluated as $\rho c_p \Delta T$ where $\Delta T = 
    \inptext{Ignition\_Source\_Magnitude} / [\#\ \textsf{RK Stages}]$.
  \item[$\inptext{Ignition\_Method}=2$]
    This value represents a percentage value that is used to add a percentage
    of the reference total energy, $(\rho e_t)_{\infty}$, as a source term to
    the energy equation. The energy source term is evaluated as
    $\inptext{Ignition\_Source\_Magnitude}\times(\rho e_t)_{\infty}$.
  \item[$\inptext{Ignition\_Method}=3$]
    This value represents a reaction temperature (units of \unitTemp{}) that
    is used to compute the species reaction source terms instead of using the
    actual local temperature for the solution point.
\end{InputList}
\end{nmldescbody}

%==========================================================================

\begin{nmldescitem}
Ignition_Source_Radius = 0.0  ! [m]
\end{nmldescitem}
\inp*[nmlChemistry::IgnitionSourceRadius]{}[Ignition\_Source\_Radius][]

\begin{nmldescbody}
Radius of influence ($R$) for the ignition source terms,
specified in 
\key[key:GridScaledUnits]{\emph{grid scaled units}} (\thatis meters)
as described in \sref{SEC:Grid_Scaling}.
%
The behavior depends upon the value ($R$) specified:
\\
\begin{tabular}{ @{\qquad} r @{ \quad } p{0.85\linewidth} @{} }
  $<0$ & The ignition source term will immediately decay
         as $\exp (-36 r^2 / R^2 )$ from the point specified by
\inp[nmlChemistry::IgnitionSourceLocation]{Ignition\_Source\_Location}.\\
  $=0$ & The ignition source is disabled.\\
  $>0$ & Every point that is within the ignition source radius ($R$)
         from the \inptext{Ignition\_Source\_Location}
         will get the full value of the energy source term
         based on the input value of
\inp[nmlChemistry::IgnitionSourceMagnitude]{Ignition\_Source\_Magnitude}.
         For all points outside this radius (R),
         the ignition source term will decay as
         $\exp (-36 (r-R)^3 / R^3 )$.
         The higher exponent realizes a slower initial decay.\\
\end{tabular}
\\
\Fref{FIG:Ignition_Source_Radius} compares the decay profiles.

\begin{figure}[h]
  \begin{centering}
    \insertimage[width=3.0in]{Ignition/Ignition_Source_Radius}
    \caption{Ignition source decay.}
    \label{FIG:Ignition_Source_Radius}
  \end{centering}
\end{figure}


\end{nmldescbody}

%==========================================================================

\pagebreak[3]
\begin{nmldescitem}
Ignition_Source_Location = [0.0,0.0,0.0]  ! [m]
\end{nmldescitem}
\inp*[nmlChemistry::IgnitionSourceLocation]{}[Ignition\_Source\_Location][]%

\begin{nmldescbody}
\raggedright % Prevent equation from entering right margin.
The center of the ignition source region, specified in 
\key[key:GridScaledUnits]{\emph{grid scaled units}} (\thatis meters)
as described in \sref{SEC:Grid_Scaling}.
\end{nmldescbody}

%==========================================================================

